\documentclass[11pt,a4paper]{article}
\usepackage[utf8]{inputenc}
\usepackage{amsmath,amssymb}
\usepackage{geometry}
\usepackage{hyperref}
\usepackage{booktabs}
\usepackage{longtable}
\usepackage{graphicx}
\usepackage{xcolor}
\usepackage{bm}

\geometry{margin=2.5cm}
\hypersetup{colorlinks=true,linkcolor=blue,urlcolor=blue}

\newcommand{\vect}[1]{\bm{#1}}

\title{\textbf{Mathematics and Algorithms of the Mirror Toolkit}}
\author{Mingrui Zuo \& Lifeng Ding\\Xi'an Jiaotong-Liverpool University}
\date{Version 0.6 -- June 2026}

\begin{document}
\maketitle
\tableofcontents
\newpage

% =================================================================
\section{Introduction}
% =================================================================

This document provides a detailed mathematical description of the core
algorithms implemented in the Mirror (Chiral Toolkit) platform.  It
covers the Canonical Principal Axes (CPA) alignment, the Kabsch
algorithm, Sohncke mirror transformations, pore detection and morphology
parameters, and guest-orientation accessibility mapping.

% =================================================================
\section{Canonical Principal Axes (CPA) Alignment}
% =================================================================

CPA alignment transforms a molecule into a standard orientation by
aligning its principal axes of inertia with the Cartesian axes.  This
provides a unique, orientation-independent reference frame.

\subsection{Inertia Tensor}

For a set of $N$ atoms with Cartesian positions
$\{\vect{r}_i = (x_i, y_i, z_i)\}$ and masses $m_i$, the inertia tensor
$\vect{I}$ is a $3 \times 3$ symmetric matrix:

\begin{equation}
\vect{I} = \begin{pmatrix}
I_{xx} & I_{xy} & I_{xz} \\
I_{xy} & I_{yy} & I_{yz} \\
I_{xz} & I_{yz} & I_{zz}
\end{pmatrix},
\end{equation}

with components:

\begin{align}
I_{xx} &= \sum_i m_i (y_i^2 + z_i^2), \\
I_{yy} &= \sum_i m_i (x_i^2 + z_i^2), \\
I_{zz} &= \sum_i m_i (x_i^2 + y_i^2), \\
I_{xy} &= -\sum_i m_i x_i y_i, \\
I_{xz} &= -\sum_i m_i x_i z_i, \\
I_{yz} &= -\sum_i m_i y_i z_i.
\end{align}

The coordinates are first centred at the centre of mass:

\begin{equation}
\vect{r}'_i = \vect{r}_i - \vect{r}_{\text{COM}},
\qquad
\vect{r}_{\text{COM}} = \frac{\sum_i m_i \vect{r}_i}{\sum_i m_i}.
\end{equation}

\subsection{Eigendecomposition and Principal Axes}

The inertia tensor is diagonalised via eigenvalue decomposition:

\begin{equation}
\vect{I} \vect{v}_k = \lambda_k \vect{v}_k,
\qquad k = 1, 2, 3,
\end{equation}

where $\lambda_k$ are the principal moments of inertia (eigenvalues) and
$\vect{v}_k$ are the principal axes (eigenvectors).  By convention, the
eigenvalues are sorted in ascending order: $\lambda_1 \le \lambda_2 \le
\lambda_3$.  The eigenvector $\vect{v}_1$ (corresponding to the smallest
moment) is assigned to the $+x$ axis, $\vect{v}_2$ to $+y$, and the
third axis is obtained from the cross product $\vect{v}_3 =
\vect{v}_1 \times \vect{v}_2$ to ensure a right-handed coordinate
system.

\subsection{Axis Orientation Convention}

The sign of each principal axis is ambiguous.  Mirror resolves this by
orienting each axis toward the side with greater atomic mass (for
mass-weighted mode) or more atoms (for geometric mode):

\begin{equation}
\vect{v}_k \leftarrow \begin{cases}
+\vect{v}_k & \text{if } \sum_{(\vect{r}'_i \cdot \vect{v}_k) > 0} m_i
               \ge \sum_{(\vect{r}'_i \cdot \vect{v}_k) < 0} m_i, \\
-\vect{v}_k & \text{otherwise}.
\end{cases}
\end{equation}

\subsection{Geometric Mode (PCA)}

In geometric mode, the covariance (scatter) matrix is used instead of the
inertia tensor:

\begin{equation}
\vect{S} = \sum_i \vect{r}'_i (\vect{r}'_i)^{\mathsf T},
\end{equation}

which is equivalent to PCA with equal weights.  The eigenvalue
decomposition and axis assignment proceed identically, but the
degeneracy criteria differ (see below).

\subsection{Degeneracy Handling}

A molecule is flagged as degenerate when its principal axes are not
uniquely defined.  The conditions depend on the mode:

\begin{itemize}
    \item \textbf{Inertia tensor (mass mode)}: degenerate if
          $\lambda_2 < \varepsilon$ (0D or 1D molecule).  Planar
          molecules are NOT degenerate because the perpendicular-axis
          theorem guarantees a non-zero out-of-plane moment.
    \item \textbf{Covariance matrix (geom mode)}: degenerate if any
          eigenvalue is zero (rank-deficient).  Planar molecules are
          degenerate here.
\end{itemize}

Additionally, if the eigenvalues are nearly equal (spherical or symmetric
top), the axes are considered degenerate:

\begin{equation}
\min\!\left(
\frac{|\lambda_3 - \lambda_2|}{\lambda_3},
\frac{|\lambda_2 - \lambda_1|}{\lambda_3},
\frac{|\lambda_3 - \lambda_1|}{\lambda_3}
\right) < \varepsilon_{\text{gap}} = 10^{-6}.
\end{equation}

In degenerate cases, the centred coordinates are returned without
rotation, and the rotation matrix is set to the identity.

% =================================================================
\section{Relationship to the Kabsch Algorithm}
% =================================================================

The Kabsch algorithm finds the optimal rotation matrix $\vect{R}$ that
minimises the root-mean-square deviation (RMSD) between two sets of
points $\{\vect{p}_i\}$ and $\{\vect{q}_i\}$:

\begin{equation}
\min_{\vect{R}} \sum_i \|\vect{R} \vect{p}_i - \vect{q}_i\|^2.
\end{equation}

The solution involves the cross-covariance matrix
$\vect{H} = \vect{P}^{\mathsf T} \vect{Q}$ and its singular value
decomposition (SVD):

\begin{equation}
\vect{H} = \vect{U} \vect{\Sigma} \vect{V}^{\mathsf T},
\qquad
\vect{R} = \vect{V} \vect{U}^{\mathsf T}.
\end{equation}

\paragraph{Relation to CPA.}  CPA differs from Kabsch in both purpose
and method:

\begin{itemize}
    \item \textbf{Purpose}: Kabsch aligns two different conformations of
          the same molecule.  CPA aligns a single molecule to a canonical
          reference frame independent of any external target.
    \item \textbf{Matrix}: Kabsch uses the cross-covariance between two
          point sets.  CPA uses the inertia tensor (mass-weighted
          auto-covariance) of a single point set.
    \item \textbf{Output}: Kabsch gives the rotation between two
          structures.  CPA gives the absolute orientation relative to the
          principal axes.
\end{itemize}

\paragraph{Unified view.}  Both algorithms ultimately compute an
eigendecomposition (CPA) or SVD (Kabsch) of a $3 \times 3$ matrix and
extract a rotation.  If one performs CPA on a structure and then applies
the inverse rotation, the principal axes align with Cartesian axes ---
analogous to rotating one structure onto another's principal frame.

% =================================================================
\section{Sohncke Mirror Transformation}
% =================================================================

\subsection{Reflection Matrices}

A mirror operation about a plane through the origin (or the molecular
centroid) is represented by a Householder reflection matrix.  The three
available planes correspond to:

\begin{align}
\vect{M}_{YZ} &= \operatorname{diag}(-1, 1, 1) \quad \text{(invert $x$)}, \\
\vect{M}_{XZ} &= \operatorname{diag}(1, -1, 1) \quad \text{(invert $y$)}, \\
\vect{M}_{XY} &= \operatorname{diag}(1, 1, -1) \quad \text{(invert $z$)}.
\end{align}

\subsection{Relation to Chirality}

If a molecule is chiral, its mirror image is a non-superimposable
enantiomer.  The CPA alignment of the mirrored structure (MCPA) enables
direct comparison: for a chiral molecule, MCPA will differ from CPA by
more than a simple rotation.  The difference in principal moments between
CPA and MCPA provides a quantitative measure of molecular chirality.

\subsection{Application to Sohncke Space Groups}

In crystallography, a Sohncke space group contains only orientation-
preserving operations (rotations, translations).  Mirror operations are
present only in non-Sohncke groups.  The Mirror toolkit classifies space
groups into three types:

\begin{itemize}
    \item \textbf{sgType 1} (Sohncke, enantiomorphic): space groups
          that allow chiral crystal structures.
    \item \textbf{sgType 2} (Sohncke, non-enantiomorphic): Sohncke
          groups without enantiomorphic pairs.
    \item \textbf{sgType 3} (non-Sohncke): contain improper rotations
          (mirrors, inversions).
\end{itemize}

The set of Sohncke space groups (65 groups) is hardcoded from the
International Tables of Crystallography, and the enantiomorphic subset
(11 pairs) is similarly defined.

% =================================================================
\section{Pore Detection}
% =================================================================

\subsection{Grid-based Distance Field}

The pore detection pipeline operates on a supercell of the CIF.  Atoms
are projected onto the $ab$-plane, and a dense 2D grid is created:

\begin{equation}
\vect{G} = \{(x_i, y_j) \mid x_i = x_{\min} + i\Delta x,\;
y_j = y_{\min} + j\Delta y\},
\end{equation}

with resolution $n_{\text{grid}}$ (default 200).  The distance from each
grid point to the nearest atom centre (minus that atom's VDW radius) is
computed:

\begin{equation}
d(\vect{g}) = \min_k \left( \|\vect{g} - \vect{a}_k\| - r_k^{\text{VDW}} \right),
\end{equation}

where $\vect{a}_k$ and $r_k^{\text{VDW}}$ are the position and VDW
radius of atom $k$.

\subsection{Pore Centre and Radius}

The largest connected region where $d(\vect{g}) \ge r_{\text{probe}}$
(default 1.2\,\AA) is identified via breadth-first search.  The pore
centre is the grid point within this region with the maximum $d(\vect{g})$,
and the pore radius $R_1$ is $d$ at that point.

\subsection{Cylindrical Projection}

VDW-surface points are generated via Fibonacci sphere sampling ($N$
points per atom).  These points are projected onto a cylinder of radius
$R_1$ centred at the pore centre:

\begin{align}
\theta &= \arctan2(y - y_0, x - x_0), \\
z     &= z_{\text{VDW}}.
\end{align}

The result is a $\theta$--$z$ point cloud that is binned into a 2D
histogram for visualisation.

\subsection{VDW Heatmap}

The VDW heatmap is a smooth Gaussian kernel-density estimate:

\begin{equation}
\rho(\theta, z) = \sum_p A_p \exp\!\left(
-\frac{(\theta - \theta_p)^2}{2\sigma_\theta^2}
-\frac{(z - z_p)^2}{2\sigma_z^2}
\right),
\end{equation}

where $A_p$ is the alpha weight of point $p$
($A_p = 1 - r_p / r_{\max}$, with $r_p$ the radial distance of the point
from the pore centre), and $\sigma_\theta$, $\sigma_z$ are proportional
to the VDW radius of the originating atom ($\sigma = K \cdot r^{\text{VDW}}$,
$K = 0.12$).

% =================================================================
\section{Morphology Parameters}
% =================================================================

\subsection{Circularity Index (CI)}

The angular wall-distance profile $r(\theta)$ is computed by binning VDW
points into 72 angular bins ($5^\circ$).  Within each bin, the mean
radial distance is computed.  CI is then:

\begin{equation}
\text{CI} = 1 - \frac{\sigma_r}{\mu_r},
\end{equation}

where $\mu_r$ and $\sigma_r$ are the mean and standard deviation across
bins.  CI $= 1$ indicates a perfectly circular pore; CI $< 1$ indicates
ellipticity or faceting.

\subsection{Fourier Power Spectrum}

The mean-subtracted profile $r'(\theta) = r(\theta) - \mu_r$ is
Fourier-transformed:

\begin{equation}
F_n = \sum_{k=0}^{N-1} r'(\theta_k) e^{-2\pi i n k / N},
\qquad P_n = \frac{|F_n|^2}{\sum_{k=0}^{N-1} |F_k|^2}.
\end{equation}

The pore label is assigned from the dominant frequency $n \ge 2$:
L6 ($n=6$), L4 ($n=4$), L3 ($n=3$), L2 ($n=2$).

\subsection{Corrugation Factor (CF)}

The approximate VDW surface area is:

\begin{equation}
A_{\text{VDW}} \approx N_{\text{pts}} \cdot R_1 \cdot \Delta\theta \cdot
\Delta z,
\end{equation}

where $\Delta\theta = 2\pi / n_\theta$ and $\Delta z = \Delta z_{\text{total}} / n_z$.
The smooth cylinder area is $A_{\text{smooth}} = 2\pi R_1 \Delta z_{\text{total}}$.
Then:

\begin{equation}
\text{CF} = \frac{A_{\text{VDW}}}{A_{\text{smooth}}}.
\end{equation}

\subsection{Mean Curvature (curv\_H)}

The VDW density field $\rho(\theta, z)$ is smoothed with a Gaussian
filter ($\sigma = 1.0$), and the Laplacian is approximated via finite
differences:

\begin{equation}
\nabla^2 \rho \approx \frac{\partial^2 \rho}{\partial \theta^2} +
\frac{\partial^2 \rho}{\partial z^2},
\end{equation}

\begin{equation}
\text{curv\_H} = \langle |\nabla^2 \rho| \rangle.
\end{equation}

\subsection{Helicity}

The angular profiles of adjacent $z$-layers are cross-correlated to find
the angular shift $\Delta\theta$ that maximises their alignment:

\begin{equation}
C_k(\Delta\theta) = \sum_i r_k(\theta_i) \cdot r_{k+1}(\theta_i + \Delta\theta).
\end{equation}

The helical twist angle is:

\begin{equation}
\alpha = \arctan\!\left(
\frac{\Delta\theta_{\max} \cdot R_1}{\Delta z_{\text{layer}}}
\right).
\end{equation}

\subsection{Radial Skewness (skew\_r)}

\begin{equation}
\text{skew\_r} = \frac{\langle (r - \mu_r)^3 \rangle}{\sigma_r^3}.
\end{equation}

Positive skewness indicates that the pore wall has more outward-
protruding features than inward ones.

% =================================================================
\section{Guest Orientation Scanning}
% =================================================================

\subsection{Rotation Parametrisation}

Each guest orientation is parametrised by two spherical angles
$(\theta, \phi)$ defining the direction of the molecule's $+z$ axis.
The rotation matrix that maps the $+z$ axis to direction
$(\theta, \phi)$ is computed via the Rodrigues rotation formula:

\begin{equation}
\vect{R}(\theta, \phi) = \vect{I} + \sin\alpha \,[\vect{k}]_\times +
(1 - \cos\alpha) (\vect{k} \vect{k}^{\mathsf T} - \vect{I}),
\end{equation}

where $\vect{k}$ is the unit rotation axis (cross product of source and
target directions) and $\alpha$ is the rotation angle.

\subsection{VDW Clash Detection}

For each orientation, the guest atom positions $\{\vect{g}_j\}$ are
computed from the rotation matrix and the guest centroid position
$\vect{p}$:

\begin{equation}
\vect{g}_j = \vect{R} \vect{g}_j^{(0)} + \vect{p}.
\end{equation}

A clash is registered if any guest--host atom pair violates the VDW
separation criterion:

\begin{equation}
\|\vect{g}_j - \vect{h}_i\| < s \cdot (r_j^{\text{VDW}} + r_i^{\text{VDW}}),
\end{equation}

where $s$ is the clash scale factor (default 0.75).  The fraction of
clash-free orientations is reported as the accessible fraction.

\subsection{Three-Region Partition (Rolling Sphere)}

The guest's minimum enclosing VDW sphere radius $R_{\text{env}}$ is
computed as:

\begin{equation}
R_{\text{env}} = \max_j \left( \|\vect{g}_j - \vect{g}_{\text{COM}}\| +
r_j^{\text{VDW}} \right).
\end{equation}

A rolling-sphere algorithm then determines, for each position
$(x, y)$ in the pore cross-section:

\begin{itemize}
    \item \textbf{Free} (green): the envelope sphere
          $(R = R_{\text{env}})$ fits without clashes.
    \item \textbf{Restricted} (amber): only the centroid
          $(R = 0)$ fits.
    \item \textbf{Blocked} (red): neither fits.
\end{itemize}

% =================================================================
\section{Space Group Analysis}
% =================================================================

Space-group detection uses pymatgen's \texttt{SpacegroupAnalyzer} at
multiple tolerance levels:

\begin{equation}
\text{symprec} \in \{0.01, 0.05, 0.1, 0.2\}.
\end{equation}

The confidence score is the fraction of tolerances that agree on the
assigned group.  Sohncke and enantiomorphic classifications are
determined by lookup tables compiled from the International Tables.

% =================================================================
\section{Summary of Parameter Meanings}
% =================================================================

\begin{longtable}{llp{7cm}}
\toprule
Parameter & Symbol & Physical meaning \\
\midrule
\endhead
CI & $\text{CI}$ & Pore circularity; 1 = perfect circle \\
CF & $\text{CF}$ & Wall corrugation; >1 = rough wall \\
curv\_H & $\langle|\nabla^2\rho|\rangle$ & Mean wall curvature \\
power\_2 & $P_2$ & 2-fold Fourier power (ellipticity) \\
power\_3 & $P_3$ & 3-fold Fourier power (trigonal) \\
power\_4 & $P_4$ & 4-fold Fourier power (square) \\
power\_6 & $P_6$ & 6-fold Fourier power (hexagonal) \\
helicity & $\alpha$ & Helical twist between layers (deg) \\
skew\_r & $\gamma_r$ & Wall-distance skewness \\
$R_1$ / $R_p$ & $R_1$ & Pore radius (\AA) \\
$R_e$ & $R_e$ & Expand radius (\AA) \\
PLD & $d_{\text{PLD}}$ & Pore-limiting diameter (\AA) \\
LCD & $d_{\text{LCD}}$ & Largest cavity diameter (\AA) \\
$I_1/I_3$ & $I_1/I_3$ & Principal-moment ratio \\
$\theta_{\text{steps}}$ & $N_\theta$ & Polar sampling for orientation scan \\
$\phi_{\text{steps}}$ & $N_\phi$ & Azimuthal sampling for orientation scan \\
$s_{\text{clash}}$ & $s$ & VDW clash scale factor \\
$r_{\text{probe}}$ & $r_p$ & Probe radius for pore detection (\AA) \\
$n_{\text{grid}}$ & $N_{\text{grid}}$ & Grid resolution for distance field \\
K & $K$ & Gaussian sigma scaling factor (VDW heatmap) \\
\bottomrule
\end{longtable}

\end{document}
